\documentclass[10pt,a4paper]{amsart}
\usepackage[margin=19mm]{geometry}
\usepackage{amsmath,amssymb,mathtools}
\usepackage[scr=rsfs,cal=euler]{mathalfa}
\usepackage{xfrac}
\usepackage[unicode,hidelinks,bookmarks=false]{hyperref}
\newcommand{\ie}{i.e.\ }
\newcommand{\cf}{cf.\ }
\newcommand{\resp}{respectively\ }
\newcommand{\Subsection}{\subsection}
\providecommand{\nobreakdash}{\nobreak\hbox{-}}
\setlength{\emergencystretch}{2em}
\multlinegap0pt
\usepackage{bm}
\usepackage{bbm}
\usepackage{dsfont}
\usepackage{xspace}
\usepackage{cancel}
\usepackage{mathtools}
\usepackage{amsthm}
\makeatletter
\@ifundefined{theorem}{\newtheorem{theorem}{Theorem}}{}
\@ifundefined{lemma}{\newtheorem{lemma}[theorem]{Lemma}}{}
\makeatother
\newcommand{\N}{\mathbb{N}}
\newcommand{\R}{\mathbb{R}}
\newcommand{\cD}{\mathcal{D}}
\newcommand{\dd}{\mathrm{d}}
\newcommand{\eps}{\varepsilon}
\renewcommand{\H}{\mathcal{H}}
\title[Laplacian eigenvalues for large negative Robin parameters on]{Laplacian eigenvalues for large negative Robin parameters on domains with outward peaks (Lemma 15)}
\author{Konstantin Pankrashkin, Firoj Sk, Marco Vogel}
\date{Source: arXiv:2504.04996v2 (CC BY 4.0)}
\begin{document}
\maketitle

\noindent\emph{Source and licence.} Laplacian eigenvalues for large negative Robin parameters on domains with outward peaks. Version arXiv:2504.04996v2 (CC BY 4.0), distributed under the Creative Commons Attribution 4.0 International licence, as recorded in the arXiv licence metadata for this exact version and shown on the abstract page https://arxiv.org/abs/2504.04996v2. Full rights notice: licenses/arxiv-2504.04996-CC-BY-4.0.txt.\par\smallskip

\noindent\textit{Editorial scope.} Counted unit: the Lemma printed as Lemma 15 of the article (source0.tex lines 1195--1204, source label \texttt{lem43}), delivered together with the article's own complete original proof of it (source0.tex lines 1206--1420, one \texttt{proof} environment of 215 lines). No mathematics is written, repaired or re-derived: statement and proof are the article's own text line for line. Required context retained uncounted: basic properties (lines 316--353); hhh (lines 960--963); hplus1b (lines 1168--1176); eq-uvw (lines 1179--1184); psv (lines 1187--1190). Every definition, display and label of those lines is retained unchanged. Omitted from this package and not redistributed: 1--315, 354--959, 964--1167, 1177--1178, 1185--1186, 1191--1194, 1205--1205, 1421--2020. Nothing inside a retained excerpt is omitted or altered: every retained body is delivered exactly as the article's lines are written, so no omission receipt of any kind accompanies this package. Citations: the retained excerpts carry no citation command at all, so no bibliography entry of the article is transcribed and no reference is redistributed. Reference adaptation: none was needed and none was made; every reference inside the delivered unit resolves inside it, so no unresolved reference or citation is produced. Printed slips kept literal: no printed word, symbol, hypothesis or number of the counted statement or of its proof was altered. Layout and class: the article's own class is \texttt{amsart} with packages including geometry, lmodern, amsmath, amsthm, amssymb, verbatim, esint, pdfpages, tikz, xcolor, hyperref; the wrapper is the lane's pinned \texttt{amsart} editorial wrapper at 10pt on A4, which restates the theorem-like environments the retained text needs and adds the article's own macro definitions that the retained text needs (\texttt{\textbackslash H}, \texttt{\textbackslash N}, \texttt{\textbackslash R}, \texttt{\textbackslash cD}, \texttt{\textbackslash dd}, \texttt{\textbackslash eps}), with extra wrapper packages (bm, bbm, dsfont, xspace, cancel, mathtools). The delivered numbering is the unit's own. Assets: no retained span contains a figure environment, a graphics-inclusion command, a PDF-page inclusion, a table or a TikZ picture; no image file is copied into this package, the package declares no asset dependency and no image file is redistributed with it. Licence: version arXiv:2504.04996v2 of this article is distributed under the Creative Commons Attribution 4.0 International licence, as recorded in the arXiv licence metadata for this exact version and shown on the abstract page https://arxiv.org/abs/2504.04996v2. A full rights notice is delivered at licenses/arxiv-2504.04996-CC-BY-4.0.txt. Characters: every retained excerpt is pure ASCII, so the delivered engine, XeLaTeX with its default Latin Modern text font, reports no missing character.
	\begin{lemma}\label{basic properties}
		Let $U\subset\R^m$ be a bounded Lipschitz domain,
		then the following properties hold true for the Robin Laplacian $R^U_\alpha$ (as defined in the introduction):
		\begin{itemize}
			
			\item[(i)]Scaling: For any $t>0$, $\alpha\in\R$, $j\in\N$ one has
			\[
			\lambda_j(R^{tU}_\alpha)=\dfrac{\lambda_j(R^{U}_{t\alpha})}{t^2}.
			\]
			
			\item[(ii)] For any $\alpha\in\R$, the first eigenvalue $\lambda_1(R^{U}_{\alpha})$
				is simple and the corresponding eigenfunctions have constant sign in $U$.
				
			\item[(iii)] The function
			\[
			\R\ni\alpha\mapsto\lambda_1(R^U_\alpha)\in\R
			\]
			is smooth. If the associated eigenfunction $\psi_\alpha$ is chosen positive with $\|\psi_\alpha\|_{L^2(U)}=1$, then the mapping
			\[
			\R\ni\alpha\mapsto \psi_\alpha\in L^2(U)
			\]
			is also smooth.

			\item[(iv)]
			
			 There exists $\phi\in L^{\infty}(0,\infty)$ such that
\begin{gather*}
\lambda_1(R^{U}_{\alpha})=-\dfrac{\H^{m-1}(\partial U)}{\H^m(U)}\alpha+\alpha^2\phi(\alpha)\text{ for all }\alpha>0.
\end{gather*}
			
			\item[(v)] If $N_U$ is the Neumann Laplacian in $U$, then
			\[
			\lim_{\alpha\to0}\lambda_2(R^{U}_{\alpha})=\lambda_2(N_U)>0.
			\]
			
			
		\end{itemize}
	\end{lemma}

\begin{align}
\int_0^a \int_{\omega}\Big(\partial_su &-\frac{q(d-1)}{2s}\, u\Big)^2\dd\H^{d-1}(t)\,\dd s=\int_0^a \int_{\omega}\bigg(|\partial_su|^2+\dfrac{H}{s^2}\,u^2\bigg)\dd\H^{d-1}(t)\,\dd s,\nonumber\\
H:&=\dfrac{q^2(d-1)^2-2q(d-1)}{4}\equiv \dfrac{\big(q(d-1)-1\big)^2-1}{4}.\label{hhh}
\end{align}

\begin{equation}
	\label{hplus1b}	
	\begin{aligned}
		h^{-}_{\eps,a}(u,u)&= (1-c\eps)\int_0^a \int_{\omega}\bigg(|\partial_su|^2+\dfrac{H}{s^2}\,u^2\bigg)\dd\H^{d-1}(t)\,\dd s\\
		&\quad
		+(1-c\eps)\int_0^a \dfrac{1}{\eps^2s^{2q}} r^\omega_{\eps\rho_\eps(s)}\big(u(s,\cdot),u(s,\cdot)\big)\dd s,\\
\text{with }	\rho_\eps(s)&:=\dfrac{s^q}{(1-c\eps)^2}.
	\end{aligned}
\end{equation}

\begin{align}
	u&=v+w, \label{eq-uvw}\\
	v(s,t)&:=f(s)\psi_{\eps\rho_\eps(s)}(t),\nonumber\\
	f(s)&:=\int_\omega \psi_{\eps\rho_\eps(s)}(t)u(s,t)\,\dd\H^{d-1}(t),\nonumber\\
w&:=u-v. \nonumber
\end{align}

\begin{align}
	\partial_s v:&\ (s,t)\mapsto f'(s)\psi_{\eps \rho_\eps(s)}(t)+f(s)\dfrac{\dd \psi_{\eps\rho_\eps(s)}}{\dd s}(t), \label{psv}\\
	\nabla_t v:&\ (s,t)\mapsto f(s)\nabla_t \psi_{\eps \rho_\eps(s)}(t).\nonumber
\end{align}

\begin{lemma}\label{lem43}
There exist a sufficiently small $\eps_0>0$ and a sufficiently large $B>0$ such that
for any $\eps\in(0,\eps_0)$ and any $u\in \cD_a$ decomposed as in \eqref{eq-uvw} it holds that
\begin{align*}
\|u\|^2_{L^2(\Pi_a)}&=\|f\|^2_{L^2(0,a)}+\|w\|^2_{L^2(\Pi_a)},\\
\dfrac{h^-_{\eps,a}(u,u)}{1-c\eps}&\ge (1-B\eps)\int_0^a\bigg[f'(s)^2  +\bigg(\dfrac{H}{s^2}
-\dfrac{A_\omega}{\eps (1-B\eps)^2 s^q}\bigg)	f(s)^2	\bigg]\dd s
-B\|f\|^2_{L^2(0,a)}.
\end{align*}
\end{lemma}

\begin{proof}
Let us collect important properties of the decomposition \eqref{eq-uvw}.
First,
\begin{align*}
	\int_\omega v(s,t)^2\dd \H^{d-1}(t)&=f(s)^2\int_\omega \psi_{\eps\rho_\eps(s)}(t)^2\dd \H^{d-1}(t)=f(s)^2,\\
	\int_{\Pi_a}v^2\dd\H^d&=\int_0^a 	\int_\omega v(s,t)^2\dd \H^{d-1}(t)\dd s=\int_0^a f(s)^2\dd s.
\end{align*}
Furthermore, by construction we have
\begin{equation}
	\label{vw00}
\int_\omega \psi_{\eps\rho_\eps(s)}(t)w(s,t)\dd \H^{d-1}(t)=0,
\end{equation}
in particular,
\begin{equation}
	\label{eq-vwperp}
\int_\omega v(s,t)w(s,t)\,\dd \H^{d-1}(t)=f(s)\int_\omega \psi_{\eps\rho_\eps(s)}(t)w(s,t)\,\dd \H^{d-1}(t)
=0,
\end{equation}
therefore,
\begin{align*}
	\int_\omega u(s,t)^2\dd\H^{d-1}(t)&=\int_\omega v(s,t)^2\dd \H^{d-1}(t) +2 \int_\omega v(s,t)w(s,t)\dd \H^{d-1}(t)\\
	&\quad+\int_\omega w(s,t)^2\dd \H^{d-1}(t)\\
	&=\int_\omega v(s,t)^2\dd \H^{d-1}(t) +\int_\omega w(s,t)^2\dd \H^{d-1}(t)\\
	&=f(s)^2+\int_\omega w(s,t)^2\dd \H^{d-1}(t),\\
	\|u\|^2_{L^2(\Pi_a)}&=\|f\|^2_{L^2(0,a)}+\|w\|^2_{L^2(\Pi_a)}.
\end{align*}
Due to the choice of $\psi_{\eps\rho_\eps(s)}$ and the orthogonality  \eqref{eq-vwperp}, the spectral theorem for $R^\omega_{\eps,\rho_\eps(s)}$ gives
\begin{align*}
r^\omega_{\eps\rho_\eps(s)}\big(u(s,\cdot),u(s,\cdot)\big)&=
r^\omega_{\eps\rho_\eps(s)}\big(v(s,\cdot),v(s,\cdot)\big)+r^\omega_{\eps\rho_\eps(s)}\big(w(s,\cdot),w(s,\cdot)\big)\\
&\ge
\lambda_1(R^\omega_{\eps\rho_\eps(s)}) \big\| v(s,\cdot)\big\|^2_{L^2(\omega)}
+\lambda_2(R^\omega_{\eps\rho_\eps(s)}) \big\| w(s,\cdot)\big\|^2_{L^2(\omega)}\\
&=\lambda_1(R^\omega_{\eps\rho_\eps(s)}) f(s)^2
+\lambda_2(R^\omega_{\eps\rho_\eps(s)}) \big\| w(s,\cdot)\big\|^2_{L^2(\omega)}.
\end{align*}
Using these estimates in \eqref{hplus1b} one obtains
\begin{align*}
	\dfrac{h^-_{\eps,a}(u,u)}{1-c\eps}&\ge \int_0^a\int_\omega |\partial_s u|^2\dd\H^{d-1}(t)\dd s+\int_0^a\bigg(\dfrac{H}{s^2}+\dfrac{\lambda_1(R^\omega_{\eps\rho_\eps(s)})}{\eps^2 s^{2q}}\bigg)f(s)^2\dd s\\
	&\quad+\int_0^a\int_\omega\bigg(\dfrac{H}{s^2}+\dfrac{\lambda_2(R^\omega_{\eps\rho_\eps(s)})}{\eps^2 s^{2q}}\bigg) w(s,t)^2\dd\H^{d-1}(t)\dd s.
\end{align*}
We first use Lemma \ref{basic properties}(iv) to estimate with a suitable $C_1>0$:
\begin{align*}
	\dfrac{\lambda_1(R^\omega_{\eps\rho_\eps(s)})}{\eps^2 s^{2q}}&
	\ge -\dfrac{A_\omega \eps\rho_\eps(s) +C_1\eps^2\rho_\eps(s)^2}{\eps^2s^{2q}}\\
	&=-A_\omega \dfrac{\rho_\eps(s)}{s^{2q}}\cdot\dfrac{1}{\eps}-C_1\dfrac{\rho_\eps(s)^2}{s^{2q}}\ge -\dfrac{A_\omega}{\eps(1-c\eps)^2}\cdot\dfrac{1}{s^q}-\dfrac{C_1}{(1-c\eps)^4}.
\end{align*}
In addition, using Lemma \ref{basic properties}(v) we find some $C_2>0$ and, if necessary, decrease the value of $\eps_0$ to obtain
\[
\lambda_2\big(R^\omega_{\eps\rho_\eps(s)}\big)\ge C_2\text{ for all $\eps\in(0,\eps_0)$ and $s\in(0,a)$.}
\]
Then for a suitable $C_3>0$ and an adjusted value of $\eps_0$ one has
\[
\dfrac{H}{s^2}+\dfrac{\lambda_2(R^\omega_{\eps\rho_\eps(s)})}{\eps^2 s^{2q}}\ge \dfrac{C_3}{\eps^2 s^{2q}}\text{ for all $\eps\in(0,\eps_0)$ and $s\in(0,a)$,}
\]
yielding
\begin{equation}
	\label{hminus1}
\begin{aligned}
	\dfrac{h^-_{\eps,a}(u,u)}{1-c\eps}&\ge \int_0^a\int_\omega |\partial_s u|^2\dd\H^{d-1}(t)\,\dd s\\
	&\quad+\int_0^a\bigg(\dfrac{H}{s^2}
	-\dfrac{A_\omega}{\eps(1-c\eps)^2}\cdot\dfrac{1}{s^q}-\dfrac{C_1}{(1-c\eps)^4}
	\bigg)f(s)^2\dd s\\
	&\quad+\int_0^a\int_\omega \dfrac{C_3}{\eps^2 s^{2q}}
	 w(s,t)^2\dd\H^{d-1}(t)\,\dd s.
\end{aligned}
\end{equation}

Our next aim is to study the term containing $|\partial_s u|^2$ on the right-hand side of \eqref{hminus1}. By denoting
\begin{align*}
	I_1&:=\int_\omega |\partial_s v|^2\dd\H^{d-1}(t),\\
	I_2&:=2\int_\omega \partial_s v\cdot\partial_s w \,\dd\H^{d-1}(t),\\
	I_3&:=\int_\omega |\partial_s w|^2\dd\H^{d-1}(t),
\end{align*}
we arrive at the decomposition
\begin{equation}
	\label{udecomp}
	\int_\omega |\partial_s u|^2\dd\H^{d-1}(t)=I_1+I_2+I_3.
\end{equation}

The term $I_1$ is analyzed in a straightforward way:
\begin{align*}
	I_1&=
	\int_\omega \Big( f'(s)\psi_{\eps \rho_\eps(s)}(t)+f(s)\dfrac{\dd \psi_{\eps\rho_\eps(s)}}{\dd s}(t)\Big)^2\dd\H^{d-1}(t)\\
	&=f'(s)^2+2f(s)f'(s)\int_\omega \psi_{\eps \rho_\eps(s)}(t)\dfrac{\dd \psi_{\eps\rho_\eps(s)}}{\dd s}(t)\dd\H^{d-1}(t)\\
	&\qquad+f(s)^2\int_\omega \Big|\dfrac{\dd \psi_{\eps\rho_\eps(s)}}{\dd s}(t)\Big|^2\dd\H^{d-1}(t),
\end{align*}
and using 
\[
\int_\omega \psi_{\eps \rho_\eps(s)}(t)\dfrac{\dd \psi_{\eps\rho_\eps(s)}}{\dd s}(t)\,\dd\H^{d-1}(t)
=\dfrac{1}{2}\dfrac{\dd}{\dd s}\int_\omega \psi_{\eps \rho_\eps(s)}(t)^2\dd\H^{d-1}(t)
=\dfrac{1}{2}\dfrac{\dd}{\dd s} 1=0
\]
to eliminate the middle term on the right-hand side, and by noting that last term is nonnegative we obtain
\begin{equation}
	\label{eq-i1}
	I_1\ge f'(s)^2.
\end{equation}	

As a consequence of \eqref{psv}, we have
\begin{align*}
	I_2&=I'_2+I''_2,\\
	I'_2&=2f'(s)\int_\omega \psi_{\eps \rho_\eps(s)}(t)\partial_s w(s,t)\,\dd\H^{d-1}(t),\\
	I''_2&=2f(s)\int_\omega \dfrac{\dd\psi_{\eps \rho_\eps(s)}}{\dd s}(t)\partial_s w(s,t)\,\dd\H^{d-1}(t).
\end{align*}
Using the orthogonality \eqref{vw00} we obtain
\begin{align*}
0&=\dfrac{\dd }{\dd s} \int_\omega \psi_{\eps \rho_\eps(s)}(t) w(s,t)\,\dd\H^{d-1}(t)\\
&=\int_\omega \dfrac{\dd\psi_{\eps \rho_\eps(s)}}{\dd s}(t) w(s,t)\,\dd\H^{d-1}(t)
+\int_\omega \psi_{\eps \rho_\eps(s)}\partial_s w(s,t)\,\dd\H^{d-1}(t),
\end{align*}
therefore,
\[
\int_\omega \psi_{\eps \rho_\eps(s)}(t)\partial_s w(s,t)\,\dd\H^{d-1}(t)=-\int_\omega 
\dfrac{\dd\psi_{\eps \rho_\eps(s)}}{\dd s}(t)  w(s,t)\,\dd\H^{d-1}(t).
\]
This allows one to estimate $I'_2$ by
\begin{align*}
	|I'_2|&=\bigg|\int_\omega 2f'(s)\dfrac{\dd\psi_{\eps \rho_\eps(s)}}{\dd s}(t) w(s,t)\dd\H^{d-1}(t)\bigg|\\
	&\le
	W_\eps(s) f'(s)^2 +\int_\omega w(s,t)^2\dd\H^{d-1}(t)\\
	\text{with }W_\eps(s)&:=\int_\omega\bigg|\dfrac{\dd\psi_{\eps \rho_\eps(s)}}{\dd s}(t)\bigg|^2\dd\H^{d-1}(t).
\end{align*}
Furthermore,
\begin{align*}
	|I''_2|&\le \dfrac{W_\eps(s)}{\eps}\, f(s)^2+ \eps \int_\omega \partial_s w(s,t)^2\dd\H^{d-1}(t).
\end{align*}	
We have 
\begin{equation*}
	\begin{split}
		W_\eps(s)
		&=\int_{\omega}\bigg(\dfrac{\dd \psi_\sigma}{\dd\sigma}(t)\Big|_{\sigma=\eps\rho_\eps(s)}
		\dfrac{\dd\big(\eps \rho_\eps(s)\big)}{\dd s}\bigg)^2\dd\H^{d-1}(t)
		\\
		&=\frac{q^2s^{2q-2}}{(1-c\eps)^4}\,\eps^2\bigg\|\dfrac{\dd \psi_\sigma}{\dd\sigma}\Big|_{\sigma=\eps\rho_\eps(s)}\bigg\|^2_{L^2(\omega)},
	\end{split}
\end{equation*}
and by Lemma \ref{basic properties}(iii) we can find $C_4>0$ such that
\[
\bigg\|\dfrac{\dd \psi_\sigma}{\dd\sigma}\Big|_{\sigma=\eps\rho_\eps(s)}\bigg\|^2_{L^2(\omega)}\le C_4 \quad\text{for all $s\in(0,a)$ and $\eps\in(0,\eps_0)$,}
\]
then for some $C_5>0$ it holds
\[
W_\eps(s)\le C_5\eps^2\text{ for all $s\in(0,a)$ and $\eps\in(0,\eps_0)$.}
\]
Then
\begin{align*}
	|I'_2|&\le C_5\eps^2 f'(s)^2+\big\|w(s,\cdot)\big\|^2_{L^2(\omega)},\\
	|I''_2|&\le C_5\eps f(s)^2+\eps\big\|\partial_s w(s,\cdot)\big\|^2_{L^2(\omega)},
\end{align*}
and
\begin{align*}
	I_2&\ge -|I'_2|-|I''_2|\ge C_5\eps^2 f'(s)^2- C_5\eps f(s)^2-\eps\big\|\partial_s w(s,\cdot)\big\|^2_{L^2(\omega)}
	-\big\|w(s,\cdot)\big\|^2_{L^2(\omega)}.
\end{align*}
By using the last inequality and \eqref{eq-i1} in \eqref{udecomp} we arrive at
\begin{align*}
	\big\|\partial_s u(s,\cdot)\big\|^2_{L^2(\omega)}&\ge (1-C_5\eps^2)f'(s)^2-C_5\eps f(s)^2\\
	&\qquad +(1-\eps)\big\|\partial_s w(s,\cdot)\big\|^2_{L^2(\omega)}-\big\|w(s,\cdot)\big\|^2_{L^2(\omega)}.
\end{align*}
By using the last inequality to estimate the first summand on the right-hand side of \eqref{hminus1}
we obtain
\begin{align*}
		\dfrac{h^-_{\eps,a}(u,u)}{1-c\eps}&\ge 
		\int_0^a\bigg[(1-C_5\eps^2)f'(s)^2+\bigg(\dfrac{H}{s^2}
		-\dfrac{A_\omega}{\eps(1-c\eps)^2 s^q}-\dfrac{C_1}{(1-c\eps)^4}-C_5\eps\bigg)
		f(s)^2	\bigg]\dd s\\
		&\quad+\int_0^a\int_\omega \bigg[(1-\eps)\partial_s w(s,t)^2+\bigg(\dfrac{C_3}{\eps^2 s^{2q}}
		-1\bigg)w(s,t)^2\bigg]\dd\H^{d-1}(t)\,\dd s.
\end{align*}
We additionally adjust $\eps_0$ such that the expression in the last integral
becomes non-negative for all $s\in(0,a)$ and all $\eps\in(0,\eps_0)$ and choose $C_6>0$
to have
\[
\dfrac{C_1}{(1-c\eps)^4}+C_5\eps \le C_6 \text{ for all }\eps\in(0,\eps_0),
\]
which results in
\begin{equation}
	\label{hminus3}
	\dfrac{h^-_{\eps,a}(u,u)}{1-c\eps}\ge \int_0^a\bigg[(1-C_5\eps^2)f'(s)^2 
	+\bigg(\dfrac{H}{s^2}
	-\dfrac{A_\omega}{\eps(1-c\eps)^2 s^q}-C_6\bigg)
	f(s)^2	\bigg]\dd s.
\end{equation}
Recall that due to the explicit expression \eqref{hhh} for $H$ and the one-dimensional Hardy inequality we have
\begin{align*}
\dfrac{\big(q(d-1)-1\big)^2}{4}\int_0^a \dfrac{1}{s^2}\, f(s)^2\dd s
&\le \int_0^a \bigg(f'(s)^2+\dfrac{\big(q(d-1)-1\big)^2-1}{4 s^2}f(s)^2\bigg)\dd s\\
&\equiv \int_0^a \bigg(f'(s)^2+\dfrac{H}{s^2}f(s)^2\bigg)\dd s,
\end{align*}
yielding
\[
\int_0^a \dfrac{1}{s^2}\, f(s)^2\dd s\le C_7 \int_0^a \bigg(f'(s)^2+\dfrac{H}{s^2}f(s)^2\bigg)\dd s,
\quad C_7:=\dfrac{4}{\big(q(d-1)-1\big)^2}.
\]
This implies
\begin{align*}
\int_0^a&\bigg[(1-C_5\eps^2)f'(s)^2 +\dfrac{H}{s^2} f(s)^2\bigg]\dd s\\
&=(1-C_5\eps^2)\int_0^a\bigg(f'(s)^2+\dfrac{H}{s^2}f(s)^2\bigg)\dd s +H C_5 \eps^2 \int_0^a \dfrac{1}{s^2}\, f(s)^2\dd s\\
&\ge (1-C_5\eps^2)\int_0^a\bigg(f'(s)^2+\dfrac{H}{s^2}f(s)^2\bigg)\dd s- |H|C_5C_7\eps^2 \int_0^a \bigg(f'(s)^2+\dfrac{H}{s^2}f(s)^2\bigg)\dd s\\
&= (1-C_8\eps^2)\int_0^a\bigg(f'(s)^2+\dfrac{H}{s^2}f(s)^2\bigg)\dd s\quad
\text{ with }C_8:=C_5+|H|C_5C_7.
\end{align*}
By using this inequality on the right-hand side of \eqref{hminus3} we arrive at
%kp detailed the computation below
\begin{align*}
	\dfrac{h^-_{\eps,a}(u,u)}{1-c\eps}&\ge (1-C_8\eps^2)\int_0^a\bigg(f'(s)^2 +\dfrac{H}{s^2}f(s)^2\bigg)\,\dd s-\int_0^a\dfrac{A_\omega}{\eps s^q (1-c\eps)^2}	f(s)^2\dd s
-C_6\|f\|^2_{L^2(0,a)}.
\end{align*}
We assume additionally $\eps_0\in(0,1)$, then $\eps^2<\eps$ for all $\eps\in(0,\eps_0)$, and we obtain the claim by choosing any $B\ge \max\{C_6,C_8\}$ such that
\[
\dfrac{1}{(1-c\eps)^2}\le\dfrac{1}{1-B\eps} \text{ for all }\eps\in(0,\eps_0) \qedhere
\]
holds.
\end{proof}

\end{document}
